\documentclass{article}
\usepackage{graphicx}
\usepackage{amssymb}
\usepackage{amsmath}
\usepackage{amsthm}
\usepackage[%  
    colorlinks=true,
    pdfborder={0 0 0},
    linkcolor=red
]{hyperref}

\newtheorem{theorem}{Theorem}
\newtheorem{definition}{Definition}
\newtheorem{example}{Example}
\newtheorem{remark}{Remark}

\begin{document}

\title{Notes on Multivariable Calculus}
\author{Soumyajit De}

\maketitle
%\tableofcontents
%\newpage
\section{Geometric Intuition and Motivation}
What are we dealing with:
\begin{itemize}
\item $f(x,y)=x^2+y$, $f:\mathbb{R}^2\rightarrow \mathbb{R}$ [functional]
\item $f(x,y)=\begin{bmatrix} 
3x \\
2y 
\end{bmatrix}$, $f:\mathbb{R}^2\rightarrow \mathbb{R}^2$ [vector-valued functions]
\item $f(x,y)=\begin{bmatrix} 
x^2 \\
y^2 \\
2x-y 
\end{bmatrix}$, $f:\mathbb{R}^2\rightarrow \mathbb{R}^3$ [parametric surfaces]
\end{itemize}
Different ways to visualize:
\begin{itemize}
\item functionals can be represented as graphs with one of the axes (typically the $z$-axis) serving for the function value itself, or contour plots
\item vector-valued functions $f:\mathbb{R}^2\rightarrow \mathbb{R}^2$ can be visualized as a vector field.
\item vector-valued functions $f:\mathbb{R}^2\rightarrow \mathbb{R}^2$ can also be visualized by watching the entire input space in $\mathbb{R}^2$ move to the output space in $\mathbb{R}^2$.
\item parametric surfaces for visualizing a mapping of a plane (say, $\mathbb{R}^2$) move into a surface embedded in $\mathbb{R}^3$.
\end{itemize}

\begin{figure}
\centering
\includegraphics[width=.3\linewidth]{calc3/001.png}\quad\includegraphics[width=.3\linewidth]{calc3/002.png}\quad\includegraphics[width=.3\linewidth]{calc3/003.png}
\caption{Slicing Graphs}
\label{figure:slicing}
\end{figure}

\subsubsection{Visualizing Functionals: Graphs and Contour plots}
For 3-dimensional graphs, slicing the graph for a fixed value of the input can give us some insight. For example, in the graph of $f(x,y)=\cos(x)\sin(y)$ (see fig \ref{figure:slicing}), we can set $x=0$ or $y=\frac{\pi}{2}$ and examine the behavior.

If we slice the graph at $z$ for some constant value. That would give rise to the contour plot of the graph (see fig \ref{figure:contour}).

\begin{figure}
\centering
\includegraphics[width=.3\linewidth]{calc3/005.png}\quad\includegraphics[width=.3\linewidth]{calc3/006.png}\quad\includegraphics[width=.3\linewidth]{calc3/007.png}
\caption{Contour Plots}
\label{figure:contour}
\end{figure}

\subsubsection{Visualizing parametric surfaces}
Parameter here is a fancy name for the input, e.g.
$f(t)=\begin{bmatrix} 
t\cos(t) \\
t\sin(t) 
\end{bmatrix}$.\\
Best way to visualizing these objects is to look at just the output space of the function (see fig \ref{figure:parametric_curve}). One thing to keep in mind is that, there can be multiple functions which look the same way in the output space, if we don`t keep track of the
\begin{figure}
\centering
\includegraphics[width=.7\linewidth]{calc3/008.png}
\caption{Parametric surface}
\label{figure:parametric_curve}
\end{figure}
\end{document}